% Part: methods
% Chapter: proofs
% Section: introduction

\documentclass[../../../include/open-logic-section]{subfiles}

\begin{document}

\olfileid[pt]{mth}{prf}{int}

\olsection{Introdução}

Com base em suas experiências na lógica introdutória, você talvez se
sinta à vontade com um sistema de !!{derivation}---provavelmente um
sistema de !!{derivation} de dedução natural ou no estilo Fitch, ou
talvez um sistema de árvores de prova. Você provavelmente se lembra de
fazer provas nesses sistemas, seja provando !!a{formula}, seja
mostrando que um dado argumento é válido. Para fazer isso, você
aplicava as regras do sistema até obter o resultado final desejado. Ao
raciocinar \emph{sobre} lógica, também provamos coisas, mas na maioria
dos casos não estamos usando um sistema de !!{derivation}. De fato, a
maior parte das provas que consideramos é feita em português (talvez
com alguma linguagem simbólica inserida) e não inteiramente na
linguagem da lógica de primeira ordem.  Ao construir tais provas, você
pode a princípio ficar perdido---como provo algo sem um sistema de
!!{derivation}?  Como começo? Como sei se minha prova está correta?

Antes de tentar uma prova, é importante saber o que é uma prova e
como construir uma.  Como o nome sugere, uma \emph{prova} destina-se
a mostrar que algo é verdadeiro. Você pode pensar nisso em termos de um
diálogo---alguém lhe pergunta se algo é verdadeiro, digamos, se todo
primo diferente de dois é um número ímpar. Responder ``sim'' não é
suficiente; a pessoa pode querer saber \emph{por quê}. Nesse caso,
você lhe daria uma prova.

No discurso cotidiano, pode ser suficiente esboçar uma resposta, ou
dar uma resposta incompleta. Na lógica e na matemática, porém,
queremos uma prova rigorosa---queremos mostrar que algo é verdadeiro
além de \emph{qualquer} dúvida. Isso significa que cada passo de nossa
prova deve ser justificado, e a justificativa deve ser convincente
(isto é, a suposição que você está usando é de fato assumida no
enunciado do teorema que você está provando, as definições que você
aplica devem ser aplicadas corretamente, as justificativas invocadas
devem ser inferências corretas, etc.).

Em geral, estamos provando algum enunciado. Chamamos os enunciados que
provamos por vários nomes: proposições, teoremas, lemas ou corolários.
Uma proposição é um enunciado básico digno de prova: importante o
suficiente para ser registrado, mas talvez não particularmente
profundo nem aplicado com frequência. Um teorema é uma proposição
significativa e importante. Sua prova é frequentemente dividida em
vários passos, e às vezes recebe o nome da pessoa que primeiro o provou
(por exemplo, o Teorema de Cantor, o teorema de
L\"owenheim--Skolem) ou do fato de que trata (por exemplo, o teorema da
completude).  Um lema é uma proposição ou teorema usado na prova de um
resultado mais importante. De modo confuso, às vezes os lemas são
resultados importantes em si mesmos, e também recebem o nome da pessoa
que os introduziu (por exemplo, o Lema de Zorn). Um corolário é um
resultado que decorre facilmente de outro.

Um enunciado a ser provado frequentemente contém suposições que
esclarecem sobre quais tipos de coisas estamos provando algo. Ele pode
começar com ``Seja $!A$ !!a{formula} da forma $!B \lif !C$'' ou
``Suponha $\Gamma \Proves !A$'' ou algo do gênero.  Essas são as
\emph{hipóteses} da proposição, teorema ou lema, e você pode assumi-las
como verdadeiras em sua prova. Elas restringem o que estamos provando e
também introduzem alguns nomes para os objetos sobre os quais estamos
falando. Por exemplo, se sua proposição começa com ``Seja $!A$
!!a{formula} da forma $!B \lif !C$'', você está provando algo apenas
sobre todas as fórmulas de certo tipo (a saber, condicionais), e fica
entendido que $!B \lif !C$ é um condicional arbitrário sobre o qual sua
prova falará.

\end{document}
